%%%%%%%%%%%%%%%%%%%%%%%%%%%%%%%%%%%%%%%%%%%%%
\section{The Vlasov-Poisson system in spherical symmetry}
\label{Sec:VPSphSym}
%%%%%%%%%%%%%%%%%%%%%%%%%%%%%%%%%%%%%%%%%%%%%

In this section we formulate the Vlasov-Poisson system with an external potential and reduce it to spherical symmetry, first for a gas whose particles have arbitrary angular momenta and then for the reduced model in which all of them have the same magnitude of the angular momentum. Next, we introduce the dimensionless variables used in our numerical computations, the action-angle variables associated with a time-independent potential and, finally, the observables on which our analysis of the mixing phenomenon is based.

We consider a collisionless gas consisting of identical particles which are subject to a time-independent external potential $\Phi_{ext}$ and to the gravitational potential $\Phi_{gas}$ generated by the gas itself. Since all the gas particles have the same mass, it is convenient to normalize the DF such that $f(t,x,p) d^3 x\, d^3 p$ is the mass of the gas contained in the volume element $d^3 x\, d^3 p$ of phase space, where from now on $p$ denotes the momentum per unit mass.\footnote{This is equivalent to working in units for which the mass of each gas particle is one, as in Ref.~\cite{pRoS20}.} With these conventions, the Vlasov-Poisson system consists of the Vlasov equation
\begin{equation}
\frac{\partial f}{\partial t} + p\cdot\nabla_x f - (\nabla_x\Phi)\cdot\nabla_p f = 0,\qquad
\Phi := \Phi_{ext} + \Phi_{gas},
\label{Eq:Vlasov}
\end{equation}
and the Poisson equation
\begin{equation}
\Delta\Phi_{gas} = 4\pi G\rho,\qquad
\rho(t,x) := \int\limits_{\Real^3} f(t,x,p) d^3 p,
\label{Eq:Poisson}
\end{equation}
where $G$ denotes Newton's constant and $\rho$ is the mass density of the gas. To obtain a unique solution of Eq.~(\ref{Eq:Poisson}) we require that $\Phi_{gas}(t,x)\to 0$ for $|x|\to \infty$.

For the external potential we choose H\'enon's isochrone potential~\cite{mH59,jB14,BinneyTremaine-Book}
\begin{equation}
\Phi_{ext}(x) = -\frac{GM}{b + \sqrt{b^2 + r^2}},\qquad r := |x|,
\label{Eq:Isochrone}
\end{equation}
which is generated by a spherical matter distribution of total mass $M$. The parameter $b\geq 0$ is a length scale which describes the transition from a harmonic potential (when $r\ll b$) to the potential of a point particle of mass $M$ (when $r\gg b$). For $b = 0$, Eq.~(\ref{Eq:Isochrone}) reduces to the Kepler potential $\Phi_{ext} = -GM/r$, which is the external potential considered in Ref.~\cite{mHgRmScS25}. The reason for this choice is that in this potential the period of the radial motion depends only on the energy of the particle~\cite{mH59,jB14}, and that the corresponding action-angle variables are known in closed form~\cite{BinneyTremaine-Book}. When the self-gravity of the gas is neglected, this provides exact solutions against which our numerical methods can be validated.


\subsection{Reduction to spherical symmetry}
\label{SubSec:SphSym}

A spherically symmetric gas configuration is described by a DF which is invariant with respect to simultaneous rotations of $x$ and $p$. Such a DF depends on $(x,p)$ only through the radius $r$, the radial component $p_r := x\cdot p/r$ of the momentum and the magnitude $L := |x\wedge p|$ of the angular momentum, that is,
\begin{equation}
f(t,x,p) = \mathcal{F}(t,r,p_r,L).
\label{Eq:SphSymDF}
\end{equation}
In this case, the density $\rho$ and the potential $\Phi_{gas}$ depend only on $t$ and $r$, and the Vlasov equation~(\ref{Eq:Vlasov}) reduces to
\begin{equation}
\frac{\partial \mathcal{F}}{\partial t} + p_r\frac{\partial \mathcal{F}}{\partial r} - \frac{\partial V_L}{\partial r}\frac{\partial \mathcal{F}}{\partial p_r} = 0,\qquad
V_L(t,r) := \Phi(t,r) + \frac{L^2}{2r^2},
\label{Eq:VlasovSph}
\end{equation}
with $V_L$ the effective potential for the radial motion. Since no derivative with respect to $L$ appears in Eq.~(\ref{Eq:VlasovSph}), the gas particles with a given value of $L$ behave like a one-dimensional gas in the potential $V_L$. However, particles with different values of $L$ interact with each other through the potential $\Phi_{gas}$, which is determined by the spherically reduced Poisson equation
\begin{equation}
\frac{1}{r^2}\frac{\partial}{\partial r}\left( r^2\frac{\partial\Phi_{gas}}{\partial r} \right) = 4\pi G\rho,\qquad
\rho(t,r) = \frac{2\pi}{r^2}\int\limits_0^\infty\int\limits_{-\infty}^\infty \mathcal{F}(t,r,p_r,L)\, dp_r\, L dL,
\label{Eq:PoissonSph}
\end{equation}
see for instance Ref.~\cite{pDeJmAeMdN17}. The solution of Eq.~(\ref{Eq:PoissonSph}) which is regular at the center and vanishes at infinity is
\begin{equation}
\Phi_{gas}(t,r) = -4\pi G\int\limits_0^\infty \frac{\rho(t,s)}{\max\{ r,s \}} s^2 ds,
\label{Eq:PhigasSol}
\end{equation}
and its gradient is $\partial\Phi_{gas}/\partial r = G m(t,r)/r^2$, where $m(t,r) := 4\pi\int_0^r \rho(t,s) s^2 ds$ is the mass of the gas contained in the sphere of radius $r$ at time $t$.

For the following, it is useful to note that when a spherically symmetric function is integrated over phase space, the volume element can be replaced with
\begin{equation}
d^3 x\, d^3 p = 8\pi^2 dr\, dp_r\, L dL.
\label{Eq:PSVolumeForm}
\end{equation}
In particular, the total mass and the total energy of the gas are given by
\begin{eqnarray}
M_{gas} &=& 8\pi^2\int\limits_0^\infty\int\limits_{-\infty}^\infty\int\limits_0^\infty \mathcal{F}(t,r,p_r,L)\, dr\, dp_r\, L dL,
\label{Eq:Mgas}\\
\mathcal{E}_{gas} &=& 8\pi^2\int\limits_0^\infty\int\limits_{-\infty}^\infty\int\limits_0^\infty \mathcal{F}(t,r,p_r,L)\left[ \frac{p_r^2}{2} + \frac{L^2}{2r^2} + \Phi_{ext}(r) + \frac{1}{2}\Phi_{gas}(t,r) \right] dr\, dp_r\, L dL,
\label{Eq:Egas}
\end{eqnarray}
both of which are conserved in time. The conservation of $\mathcal{E}_{gas}$ provides a test for the accuracy of our numerical evolutions, see section~\ref{Sec:Validation}.


\subsection{The reduced model with fixed angular momentum}
\label{SubSec:FixedL}

A further simplification, which has been considered in Refs.~\cite{pDeJmAeMdN17,pRoS20,mHgRmScS25}, is obtained by assuming that all the gas particles have the same magnitude $L_0 > 0$ of the angular momentum.\footnote{Note that in Ref.~\cite{mHgRmScS25} the symbol $L$ refers to the square of the magnitude of the angular momentum.} In this case, the DF is of the form
\begin{equation}
\mathcal{F}(t,r,p_r,L) = \mathcal{F}_0(t,r,p_r)\delta(L - L_0),
\label{Eq:FixedL}
\end{equation}
and Eqs.~(\ref{Eq:VlasovSph},\ref{Eq:PoissonSph}) yield
\begin{equation}
\frac{\partial \mathcal{F}_0}{\partial t} + p_r\frac{\partial \mathcal{F}_0}{\partial r} - \frac{\partial V_{L_0}}{\partial r}\frac{\partial \mathcal{F}_0}{\partial p_r} = 0,\qquad
\rho(t,r) = \frac{2\pi L_0}{r^2}\int\limits_{-\infty}^\infty \mathcal{F}_0(t,r,p_r)\, dp_r,
\label{Eq:VlasovFixedL}
\end{equation}
where the potential $\Phi_{gas}$ is still given by Eq.~(\ref{Eq:PhigasSol}). Likewise, Eqs.~(\ref{Eq:Mgas},\ref{Eq:Egas}) hold with $\mathcal{F}$ replaced by $\mathcal{F}_0$, $L$ by $L_0$, and the integral over $L$ by the factor $L_0$. Equations~(\ref{Eq:VlasovFixedL},\ref{Eq:PhigasSol}) constitute a Vlasov-Poisson system on a two-dimensional phase space with coordinates $(r,p_r)$. It describes a gas of concentric spherical shells which move in the effective potential $V_{L_0}$. The centrifugal term in this potential prevents the shells from reaching the center, and the gas configurations considered in this article are confined to a region of the form $0 < r_{in}\leq r\leq r_{out}$. In spite of its simplicity, the reduced model retains the two ingredients which are relevant for this article, namely the dependency of the frequency of the radial motion on the orbit, which is responsible for the mixing, and the self-gravity of the gas.

In the second situation considered in this article, $\mathcal{F}$ is a regular function of $L$ whose support is contained in an interval $[L_1,L_2]$ with $0 < L_1 < L_2$. In this case, the phase space is three-dimensional, and the orbits with different values of $L$ have different frequencies.


\subsection{Dimensionless variables}
\label{SubSec:Units}

In our numerical computations we use dimensionless variables, which are obtained by measuring masses in units of $M$, lengths in units of $b$ and times in units of the dynamical time
\begin{equation}
t_{dyn} := \sqrt{\frac{b^3}{GM}}.
\label{Eq:tdyn}
\end{equation}
Accordingly, $p_r$ is measured in units of $\sqrt{GM/b}$, $L$ in units of $\sqrt{GMb}$, the potentials and the energy per unit mass in units of $GM/b$, and $\mathcal{F}$ in units of $M/(GMb)^{3/2}$. In terms of these variables, all the equations of this section retain their form with $G = M = b = 1$; in particular, the isochrone potential is $\Phi_{ext}(r) = -1/(1 + \sqrt{1 + r^2})$. The resulting system contains no free parameters, and the strength of the self-gravity is controlled by the ratio $M_{gas}/M$ between the mass of the gas and the mass which generates the external potential. In the limit $M_{gas}/M\to 0$ the potential $\Phi_{gas}$ can be neglected, and one recovers the model analyzed in Ref.~\cite{pRoS20}.

In the Kepler case $b = 0$ the external potential does not provide a length scale. However, in this case the system~(\ref{Eq:VlasovSph},\ref{Eq:PoissonSph}) is invariant with respect to the rescaling
\begin{equation}
(t,r,p_r,L,\mathcal{F},\Phi)\mapsto \left( \alpha^{3/2} t,\alpha r,\alpha^{-1/2} p_r,\alpha^{1/2} L,\alpha^{-3/2}\mathcal{F},\alpha^{-1}\Phi \right),\qquad \alpha > 0,
\label{Eq:Rescaling}
\end{equation}
which leaves $M_{gas}$ unchanged. Therefore, $b$ can be replaced with an arbitrary length in the definitions above, such that $\Phi_{ext}(r) = -1/r$, and in the reduced model this freedom can be used to fix the value of $L_0$.

To give an idea of the physical time scale, we note that $M = 10^5 M_\odot$ and $b = 1$~pc, which are of the order of magnitude of the mass and of the core radius of a globular cluster~\cite{BinneyTremaine-Book}, yield $t_{dyn}\approx 4.7\times 10^4$~yr. As will be seen below, $2\pi t_{dyn}$ is a lower bound for the periods of the radial motion in the isochrone potential.


\subsection{Steady states and action-angle variables}
\label{SubSec:AA}

Suppose that the total potential $\Phi$ is independent of time. Then, the energy per unit mass
\begin{equation}
E := \frac{p_r^2}{2} + V_L(r)
\label{Eq:EDef}
\end{equation}
is conserved along the trajectories of the gas particles, and it follows that any DF of the form
\begin{equation}
\mathcal{F}(r,p_r,L) = \psi(E,L),
\label{Eq:SteadyState}
\end{equation}
with a nonnegative function $\psi$, is a time-independent solution of the Vlasov equation~(\ref{Eq:VlasovSph}). In order for Eq.~(\ref{Eq:SteadyState}) to provide a solution of the Vlasov-Poisson system, the potential $\Phi_{gas}$, which enters the definition of $E$ through $V_L$, has to coincide with the one generated by the density of the DF~(\ref{Eq:SteadyState}) through Eqs.~(\ref{Eq:PoissonSph},\ref{Eq:PhigasSol}). This is a nonlinear problem for $\Phi_{gas}$, whose numerical solution is described in section~\ref{Sec:Numerics}. In the following, a solution $(\mathcal{F}_{eq},\Phi_{eq})$ of this problem, with $\Phi_{eq}$ the total potential and $\rho_{eq}$ the corresponding density, will be referred to as a self-consistent steady state.

Next, we introduce action-angle variables for the radial motion in a time-independent potential $\Phi$, following Refs.~\cite{BinneyTremaine-Book,pRoS20}. In our applications, $\Phi$ is either the external potential~(\ref{Eq:Isochrone}) or the total potential $\Phi_{eq}$ of a self-consistent steady state. In both cases $\Phi'(r) = G m_{tot}(r)/r^2$, where $m_{tot}(r)$ denotes the total mass contained in the sphere of radius $r$, including the one which generates the external potential. Since the function $r\mapsto r m_{tot}(r)$ increases strictly monotonically from $0$ to $\infty$, for each $L > 0$ the effective potential $V_L$ has a unique critical point $r = r_{min}(L)$. It is determined by the equation $r^3\Phi'(r) = L^2$, it is a non-degenerate minimum of $V_L$, and it corresponds to a circular orbit with energy $E_{min}(L) := V_L(r_{min}(L))$. Since $V_L(r)\to \infty$ for $r\to 0$ and $V_L(r)\to 0$ for $r\to \infty$, the orbits with energies in the range $E_{min}(L) < E < 0$ are bound, and they oscillate between the turning points $r_1(E,L) < r_2(E,L)$, which are the two solutions of the equation $V_L(r) = E$.

As in Ref.~\cite{pRoS20} we define the area function
\begin{equation}
A(E,L) := \oint p_r dr = 2\int\limits_{r_1(E,L)}^{r_2(E,L)} \sqrt{2(E - V_L(r))}\, dr,
\label{Eq:AreaFunction}
\end{equation}
which represents the area enclosed by the orbit in the $(r,p_r)$-plane, and whose partial derivative with respect to $E$ yields the period of the radial motion,
\begin{equation}
T(E,L) := \frac{\partial A}{\partial E}(E,L) = 2\int\limits_{r_1(E,L)}^{r_2(E,L)} \frac{dr}{\sqrt{2(E - V_L(r))}}.
\label{Eq:TDef}
\end{equation}
The action and angle variables associated with the radial motion are
\begin{equation}
J := \frac{A(E,L)}{2\pi},\qquad
Q := \frac{2\pi}{T(E,L)}\int\limits_{\gamma_r} \frac{dr}{p_r},
\label{Eq:JQDef}
\end{equation}
where $\gamma_r$ denotes the segment of the orbit in the $(r,p_r)$-plane which connects, in the clockwise direction, the inner turning point $(r_1(E,L),0)$ with the given point $(r,p_r)$. Hence, $Q/\Omega$ with $\Omega := 2\pi/T(E,L)$ is the time which has elapsed since the last passage of the particle through its pericenter; in particular, $Q = 0$ at the pericenter, $Q = \pi$ at the apocenter, and $0 < Q < \pi$ when $p_r > 0$. For each fixed value of $L$, the map $(r,p_r)\mapsto (Q,J)$ is symplectic, that is $dr\wedge dp_r = dQ\wedge dJ$, and it provides smooth coordinates on the set of bound orbits which are not circular. Since $T > 0$, the first relation in Eq.~(\ref{Eq:JQDef}) can be inverted for each $L$, which yields the energy as a function $E(J,L)$ of the action variable and the angular momentum, and the frequency of the radial motion is
\begin{equation}
\Omega(J,L) = \frac{\partial E}{\partial J}(J,L) = \frac{2\pi}{T(E,L)}.
\label{Eq:OmegaDef}
\end{equation}
\EJ{We denote the orbital frequency by $\Omega$ (it is called $\omega$ in Ref.~\cite{pRoS20}) and reserve $\omega$ for the frequency of the collective oscillation.}
In terms of the variables $(Q,J,L)$, the volume element~(\ref{Eq:PSVolumeForm}) is $8\pi^2 dQ\, dJ\, L dL$, and the Vlasov equation~(\ref{Eq:VlasovSph}) for a time-independent potential assumes the simple form
\begin{equation}
\frac{\partial \mathcal{F}}{\partial t} + \Omega(J,L)\frac{\partial \mathcal{F}}{\partial Q} = 0.
\label{Eq:VlasovAA}
\end{equation}
In particular, the steady states~(\ref{Eq:SteadyState}) are precisely those DFs which are independent of the angle variable, and they can be regarded as functions of $(J,L)$.

For the external potential~(\ref{Eq:Isochrone}) the area function can be computed explicitly, see Refs.~\cite{BinneyTremaine-Book,pRoS20}. In terms of the dimensionless variables introduced in the previous subsection, one finds
\begin{equation}
A(E,L) = 2\pi\left[ \frac{1}{\sqrt{-2E}} - c(L) \right],\qquad
E_{min}(L) = -\frac{1}{2c(L)^2},
\label{Eq:AreaIsochrone}
\end{equation}
with $c(L) := \left( L + \sqrt{L^2 + 4} \right)/2$ for the isochrone potential and $c(L) := L$ for the Kepler potential. This yields
\begin{equation}
E(J,L) = -\frac{1}{2\left[ J + c(L) \right]^2},\qquad
\Omega(J,L) = \frac{1}{\left[ J + c(L) \right]^3} = (-2E)^{3/2}.
\label{Eq:OmegaIsochrone}
\end{equation}
Therefore, the frequency depends only on the energy, it decreases monotonically when $J$ or $L$ increase, and for each fixed $L$ its maximum is attained in the limit $J\to 0$ of circular orbits. For the isochrone potential, $c(L) > 1$ implies that $\Omega < 1$, such that the periods of the radial motion are larger than $2\pi$. An explicit expression for the angle variable $Q$ in terms of $(r,p_r,L)$ can be found in section 3.5.2 of Ref.~\cite{BinneyTremaine-Book}; its inversion requires the solution of an equation of Kepler's type. In contrast to this, for a self-consistent steady state the potential $\Phi_{eq}$ is only known numerically, and in this case we compute the action-angle variables by numerical quadrature of Eqs.~(\ref{Eq:AreaFunction},\ref{Eq:TDef},\ref{Eq:JQDef}), see Appendix~\ref{App:ActionAngle}.


\subsection{Observables}
\label{SubSec:Observables}

As in Ref.~\cite{pRoS20}, we describe the macroscopic state of the gas through observables of the form
\begin{equation}
N_\varphi(t) := \int f(t,x,p)\varphi(x,p)\, d^3 x\, d^3 p,
\label{Eq:Observable}
\end{equation}
where $\varphi$ is a given test function on phase space. In this article we consider two classes of such observables.

The first one refers to the action-angle variables $(Q,J)$ belonging to a time-independent reference potential. Unless otherwise stated, this is the total potential $\Phi_{eq}$ of the self-consistent steady state whose perturbations are analyzed, which reduces to $\Phi_{ext}$ when the self-gravity of the gas is neglected. For a spherically symmetric configuration and a test function of the form $\varphi = a(Q) B(J,L)$, with real-valued functions $a$ and $B$, Eqs.~(\ref{Eq:PSVolumeForm},\ref{Eq:Observable}) and the expansion of $a$ in a Fourier series yield
\begin{equation}
N_\varphi(t) = \sum\limits_{k\in\Integer} \hat{a}_k^*\, h_k(t),\qquad
\hat{a}_k := \frac{1}{2\pi}\int\limits_0^{2\pi} a(Q) e^{-ikQ} dQ,
\label{Eq:NphiSum}
\end{equation}
with the functions
\begin{equation}
h_k(t) := 8\pi^2\int\limits_0^\infty\int\limits_0^\infty\int\limits_0^{2\pi} \mathcal{F}(t,Q,J,L) B(J,L) e^{-ikQ}\, dQ\, dJ\, L dL,\qquad k\in\Integer.
\label{Eq:hk}
\end{equation}
In the reduced model, $B$ depends only on $J$ and Eq.~(\ref{Eq:hk}) holds with $\mathcal{F}$ replaced by $\mathcal{F}_0$ and the integral over $L$ by the factor $L_0$. Since $\mathcal{F}$ is real, $h_{-k} = h_k^*$, such that it is sufficient to consider $k\geq 0$. For a steady state $h_k$ vanishes for all $k\neq 0$, whereas $h_0$ is equal to the mass of the gas weighted with the function $B$. Therefore, the functions $h_k(t)$ with $k\neq 0$ measure the departure of the gas configuration from a steady state, and their decay in time is the signature of the relaxation process we are interested in. We stress that this property relies on the use of the action-angle variables which belong to the total potential $\Phi_{eq}$. If those of the external potential are used instead, a self-consistent steady state gives rise to nonvanishing, time-independent values of $h_k$ with $k\neq 0$, as we show in section~\ref{Sec:Validation}.

The second class of observables consists of the perturbation of the gravitational potential,
\begin{equation}
\delta\Phi(t,r) := \Phi(t,r) - \Phi_{eq}(r)
 = -4\pi G\int\limits_0^\infty \frac{\rho(t,s) - \rho_{eq}(s)}{\max\{ r,s \}} s^2 ds,
\label{Eq:deltaPhi}
\end{equation}
where we have used Eq.~(\ref{Eq:PhigasSol}). For each fixed $r$ this is the difference between the values of the observable~(\ref{Eq:Observable}) with test function $\varphi(x,p) = -G/\max\{ r,|x| \}$ which belong to the DF of the gas and to the one of the steady state. In contrast to the functions $h_k(t)$, the quantity $\delta\Phi$ does not refer to any particular coordinates on phase space, and together with its gradient it is the one whose decay is addressed in Refs.~\cite{mHgRmScS25,mHmS25}. To measure its magnitude we use the root mean square
\begin{equation}
\| \delta\Phi(t) \| := \left( \frac{1}{r_b - r_a}\int\limits_{r_a}^{r_b} \delta\Phi(t,r)^2 dr \right)^{1/2}
\label{Eq:deltaPhiNorm}
\end{equation}
over an interval $[r_a,r_b]$ which covers the region occupied by the gas.
